\documentclass[conference,a4paper]{IEEEtran}
\IEEEoverridecommandlockouts

\usepackage[T1]{fontenc}
\usepackage[utf8]{inputenc}
\usepackage{cite}
\usepackage{amsmath,amssymb,amsfonts}
\usepackage{graphicx}
\usepackage{textcomp}
\usepackage{xcolor}
\usepackage{url}
\usepackage{tabularx}
\usepackage{booktabs}
\usepackage{eurosym}
\usepackage{placeins}

\newcolumntype{L}{>{\raggedright\arraybackslash}X}
\newcommand{\eur}{\ifmmode\text{\texteuro}\else\texteuro\fi}
\newcommand{\cIst}{c_{\mathrm{cur}}}
\newcommand{\cSoll}{c_{\mathrm{tgt}}}

\begin{document}

% ============================================================
%  TITEL
% ============================================================
\title{Digital Twins for Predictive Maintenance in Automotive Aftersales}

% ============================================================
%  AUTOREN  -- alle 3 Gruppenmitglieder eintragen
% ============================================================
\author{
  \IEEEauthorblockN{Maksim Bogachenkov}
  \IEEEauthorblockA{
    Syntegon \\
    Stuttgart \\
    wi23024@lehre.dhbw-stuttgart.de
  }
  \and
  \IEEEauthorblockN{Oliver Yanjie Feng}
  \IEEEauthorblockA{
    TRUMPF SE + Co. KG \\
    Ditzingen \\
    wi23030@lehre.dhbw-stuttgart.de
  }
  \and
  \IEEEauthorblockN{Joel Yerai Martinez Campo}
  \IEEEauthorblockA{
    BITBW \\
    Stuttgart \\
    wi23194@lehre.dhbw-stuttgart.de
  }
}

\maketitle

% ============================================================
%  ABSTRACT  (keine Formeln, Sonderzeichen oder Symbole!)
%  Ziel: verstaendlicher Ueberblick ueber Thema, Vorgehen,
%        wichtigste Ergebnisse und Fazit (~150-250 Woerter)
% ============================================================
\begin{abstract}
This paper investigates the effect of digital-twin-based predictive maintenance on the automotive aftersales domain, focusing on the German passenger-car market. A model-based gap analysis is conducted to compare the interval-based current state against a digital-twin-enabled target state along the full customer journey. Based on this comparison, the business potential is quantified using a transparent assumption framework anchored in published German market data. The results show that proactive, data-driven service recommendations raise the repair capture rate from 40 to 70 percent, increasing wear-part gross profit by approximately 53 EUR per vehicle per year and 263,000 EUR annually across a 5,000-vehicle fleet. Customer touchpoints grow from 1.6 to roughly 2.7 per year, each representing an additional advisory and upselling opportunity. The findings establish a transparent, contestable model of the potential that dealerships and OEMs can use to evaluate and plan the integration of digital-twin-based predictive maintenance into their aftersales operations.
\end{abstract}

\begin{IEEEkeywords}
Digital Twin, Predictive Maintenance, Connected Vehicle, Machine Learning, Automotive Aftersales
\end{IEEEkeywords}


% ============================================================
\section{Introduction}
% ============================================================

Today's vehicles are connected by default and continuously transmit operational
data \cite{lopes2025vehicledt,coppola2016connected}, so the information required
to assess their technical condition is largely available. Nevertheless, automotive aftersales remains reactive: even well-connected vehicles are serviced on fixed time or mileage
intervals rather than on the actual condition of their components \cite{national_transport_university_overview_2023}.
Wear parts---brake pads and discs, tyres, ignition components, or wipers---are typically replaced only after wear or failure, leading to unplanned repairs frequently carried out by independent third-party workshops \cite{alexander_secret_2002}. For the
dealership, which owns the customer relationship, this means losing
repair revenue and valuable interactions, reducing opportunities for
customer advice, retention, and upselling \cite{challagalla_proactive_2009,sabbagha_impact_2016}.

Digital-twin-based predictive maintenance (PdM) can turn this data into
foresight. A digital twin (DT) maintains a synchronized virtual representation of
the vehicle \cite{grieves2014}, against which machine-learning (ML) models forecast
the remaining useful life (RUL) of components and recommend replacement before
failure---supporting closer customer retention at the dealership. While the building blocks of DT-based PdM are well studied
in industrial settings \cite{wang2023review,abougrad2026xai},
their integration into automotive aftersales remains underexplored.

This paper therefore addresses the following research question: \emph{what effect does the proactive use of digital
twins and predictive maintenance have on the automotive aftersales domain?}
Focusing on the German market and working with explicitly stated assumptions,
the study (i)~clarifies the interplay of DT and PdM in aftersales, (ii)~compares the predictive target state with the interval-based current state along the full customer journey to identify where reactive servicing causes value leakage, and (iii)~derives actionable recommendations for dealerships and original equipment manufacturers (OEMs).

% ============================================================
\section{Related Work and Theoretical Background}\label{sec:related}
% ============================================================

\subsection{Digital Twins}

The concept of the digital twin was first introduced by Grieves~\cite{grieves2014} in a white paper on manufacturing excellence, where it described a virtual replica of a physical product maintained in synchrony with its real-world counterpart throughout the entire asset lifecycle. Since then, the concept has expanded into aerospace, energy, and automotive domains~\cite{wang2023review,wise_go_2000}.

In the automotive domain, this general architecture maps directly onto the connected vehicle. The physical entity is the vehicle itself, equipped with on-board sensors and an On-Board Diagnostic (OBD) interface. Its virtual representation is a cloud-hosted model mirroring component states, degradation histories, and operating conditions. The synchronising link is the over-the-air (OTA) telematics connection, which transmits telemetry data to the cloud and returns configuration updates to the vehicle~\cite{lopes2025vehicledt,coppola2016connected,zhang2009connected}. Together, this setup transforms the vehicle from a periodically inspected asset into a continuously monitored system --- a shift first documented for automotive applications by Zhang et al.~\cite{zhang2009connected}.

%\begin{figure}[h]
%  \centering
%  \includegraphics[width=0.65\linewidth]{../pictures/DT-Architecture.png}
%  \caption{Three-layer architecture of an automotive digital twin.}
%  \label{fig:dt-architecture}
%\end{figure}


%\begin{figure}[htbp]
%  \centering
%  \includegraphics[width=0.45\textwidth]{graphics}
% \caption{}
%  \label{DT-Bild}
%\end{figure}

\subsection{Vehicle Maintenance Strategies}

Vehicle maintenance strategies form a maturity hierarchy: reactive maintenance repairs only after failure \cite{abougrad2026xai}; preventive maintenance schedules service at fixed mileage or time intervals, leading to either premature replacement or undetected wear \cite{abougrad2026xai,aichi2024multilayered}; condition-based maintenance (CBM) couples the service decision to sensor thresholds \cite{aichi2024multilayered}; and PdM extends CBM by forecasting the future evolution of degradation, estimating the remaining time or distance until end-of-life is reached \cite{kumar2025suspension,abougrad2026xai}.

PdM at this scale is not a stand-alone technique but a service layer built on a DT of the vehicle. As Lopes et al.~\cite{lopes2025vehicledt} and Hadraoui et al.~\cite{aichi2024multilayered} describe, the DT supplies the synchronised virtual representation against which sensor streams are interpreted, anomalies detected, and remaining-useful-life estimates computed. In current automotive DT architectures, predictive analytics is therefore typically positioned as a dedicated microservice within the DT's analytics core \cite{lopes2025vehicledt,karthick2025simulating,abougrad2026xai}.

\subsection{Machine Learning in PdM}

Three problem classes structure machine-learning (ML)-based PdM. Classification assigns sensor patterns to discrete fault categories using Support Vector Machines, Random Forests, or Convolutional Neural Networks on time-frequency representations of vibration signals \cite{kumar2025suspension,abougrad2026xai}. Regression targets Remaining Useful Life (RUL) estimation, dominated by Long Short-Term Memory (LSTM) networks for their ability to model long-range temporal dependencies in degradation signals \cite{karthick2025simulating,abougrad2026xai}. Anomaly detection operates without explicit labels, learning nominal behaviour and flagging deviations as failure precursors \cite{karthick2025simulating,abougrad2026xai}. AbouGrad et al. report classification accuracies of approximately 89\% for vibration-based fault detection and 92\% for LSTM-based RUL estimation in automotive contexts, though they note that these figures depend strongly on dataset quality and operating conditions~\cite{abougrad2026xai}. In the aftersales context, RUL regression is the primary mechanism behind proactive service recommendations; anomaly detection serves as an early-warning layer for components without established wear curves, where labelled failure data is absent.

\subsection{Automotive Aftersales}

Aftersales encompasses post-purchase activities that help customers access supplier services and resolve product issues through continuous maintenance and upkeep \cite{mertins_servitization_1999}. In the automotive context, it functions as a dedicated business module bundling maintenance, repair, and warranty offerings to maximise product use throughout the asset lifecycle while meeting financial and customer-performance targets \cite{sabbagha_impact_2016,balinado_effect_2021}. Building post-sales relationships strengthens customer loyalty, yields actionable insights, and creates upselling and cross-selling opportunities \cite{peters1997it}, making the automotive aftersales market increasingly attractive as demand for higher-quality service expands \cite{sabbagha_impact_2016}.  

Historically, once the warranty expires, customer engagement with the dealership tends to decline. Within the automotive industry, aftersales services represent a particularly significant source of revenue, with profit margins considerably exceeding those generated through vehicle sales~\cite{macbryde_role_2006}. From an economic standpoint, the aftersales service market has been identified as substantially larger than the primary product sales market across various industries~\cite{bundschuh_how_2003,editorial_universitat_politecnica_de_valencia_impact_2015,wise_go_2000}. Retaining these customers requires better information systems and service processes~\cite{mertins_servitization_1999}.

Proactive aftersales improves customer satisfaction and retention by resolving issues before customers encounter them, enabling more productive service interactions~\cite{challagalla_proactive_2009}. This requires tailoring services to customer needs through tight integration of operational and informational data~\cite{mertins_servitization_1999}. ML and DT-based approaches enable predictive interventions and better asset lifecycle optimisation, delivering cost reduction, improved customer satisfaction, and extended asset lifetimes~\cite{gonzalez-pridaDigitalTransformationAftersales2025}.


\subsection{Synthesis and Research Gap}

The preceding subsections establish three building blocks: the DT as a synchronised vehicle model, PdM as a prognostic analytics layer built on top of it, and aftersales as the customer-facing business context in which service decisions are communicated and acted upon. Together, they form a potential end-to-end system in which the DT continuously feeds component-state data to ML-based PdM models, whose forecasts are surfaced through aftersales processes to enable proactive service appointments and retain customers at the dealership.

While the referenced contributions establish the technical foundations of DT-based PdM, existing work concentrates either on algorithmic aspects in isolation \cite{kumar2025suspension,abougrad2026xai} or on industrial applications \cite{wang2023review}. The integration of DT-based PdM into automotive aftersales business models, particularly regarding its implications for customer interaction, remains underexplored. The present work addresses this gap.
% ============================================================
\section{Methodology}
% ============================================================

Building on the theoretical background established in Section~\ref{sec:related},
this paper applies a model-based gap analysis to answer the research question.

\paragraph{Gap analysis.}
A model-based rather than experimental approach is adopted for two reasons:
deploying a prototype dealership system lies outside the scope of this paper,
and no public in-service wear-part failure dataset for the German passenger-car
market exists at the required resolution~\cite{abougrad2026xai}. Explicit
assumptions anchored to published German market statistics therefore substitute
for empirical calibration, with the trade-off that the result is a transparent,
contestable estimate of potential rather than a measured outcome.
The analysis proceeds in two stages. First, a lifecycle comparison
(Section~\ref{sec:lifecycle}) contrasts the interval-based current state against
a DT/PdM-enabled target state along the customer journey, identifying the
structural points at which reactive servicing causes value leakage. Second, a
quantitative potential assessment translates the identified gap into economic
KPIs. A single economic lever---the repair capture rate $c$, the share of
addressable wear-part work retained at the franchised dealership---is used to
quantify the gap; its target value is derived analytically from the mechanism
assumptions rather than assumed directly, ensuring internal consistency. A
one-way sensitivity analysis over $c$ confirms that the economic case holds
across a plausible parameter band~\cite{cooper_gap_2015}.


% ============================================================
\section{Gap Analysis}
% ============================================================

\subsection{Current- and Target-State Lifecycle Comparison}\label{sec:lifecycle}

\begin{figure}[htbp]
  \centering
  \includegraphics[width=1.0\columnwidth]{../pictures/Lifecycle Aftersales}
  \caption{Aftersales lifecycle: current interval-based state (left) versus DT/PdM-enabled target state (right).}
  \label{fig:lifecycle}
\end{figure}

Figure~\ref{fig:lifecycle} contrasts the two states side by side, tracing the
customer journey from vehicle operation through service execution to post-repair
follow-up, and reveals at which points reactive servicing causes value leakage
and where a predictive architecture would redirect that value back to the
franchised dealership.

\paragraph{Current state.}
In the interval-based current state, the customer journey is driven by fixed
manufacturer schedules and visible failure signals. The vehicle operates without
continuous component-state monitoring; no wear forecast is computed during
normal driving. Service appointments are initiated either by a time- or
mileage-triggered reminder displayed on the instrument cluster, or by an acute
fault manifesting as a warning lamp or perceptible degradation in braking or
ride performance. Both triggers are late-stage events: the mileage interval is
set conservatively to cover the worst-case fleet population, so most vehicles
reach the workshop with substantial remaining part life; the fault-driven trigger
often arrives after the wear limit has already been exceeded~\cite{abougrad2026xai}.
Because no proactive outreach precedes the fault, the customer exercises free
choice of workshop at the moment of highest urgency, and the majority of these
unplanned repairs are captured by independent third-party
garages~\cite{dat2026report}. This leakage intensifies with
vehicle age---franchised-dealer share drops from 86\,\% for vehicles under
three years old to just 39\,\% for those over six years~\cite{dat2026report}---and
with cost-driven customer behavior: 43\,\% of German vehicle owners now perform
only the minimum maintenance~\cite{dat2026report}. The dealership's contact with
the customer is therefore reduced to the statutory inspection cycle, limiting both
retention and upselling opportunity.

\paragraph{Target state.}
In the DT/PdM-enabled target state, the vehicle's digital twin continuously
aggregates sensor streams, OBD signals, and OTA telemetry to maintain an
up-to-date virtual representation of each monitored
component~\cite{lopes2025vehicledt}. Machine-learning models operating on this
virtual representation compute remaining-useful-life estimates for wear parts
such as brake pads and discs, tyres, and ignition
components~\cite{kumar2025suspension,karthick2025simulating}. When the forecast
horizon falls below a configurable threshold, the system generates a proactive
service recommendation communicated to the customer through the OEM's
connected-car interface or the dealership's CRM before any visible fault has
occurred~\cite{challagalla_proactive_2009,gonzalez-pridaDigitalTransformationAftersales2025}.
The appointment is scheduled at the customer's convenience, the required parts
are pre-ordered, and the workshop interaction is transformed from an unplanned,
fault-driven event into a planned, evidence-based service visit. Post-service,
the digital twin resets the component baseline and resumes monitoring, closing
the lifecycle loop. The dealership thereby intercepts wear-part replacements
that would otherwise be lost to the independent aftermarket and gains structured
contact points along the customer journey that the current state does not provide.

\subsection{Assumption Framework}

Because public, in-service failure datasets representative of this setting are
largely absent~\cite{abougrad2026xai}, the comparison is built on an explicit assumption
framework rather than on a proprietary dataset. Assumptions are divided into
\emph{context} assumptions, which fix the reference case (Table~\ref{tab:context}),
\emph{wear-part} assumptions, which size the addressable work
(Table~\ref{tab:wear}), and \emph{mechanism} assumptions, which govern the
transition from the current to the target state (Table~\ref{tab:mech}). Each is
anchored to a published figure where one exists and flagged as a set value otherwise.
 
\begin{table}[h!]
\caption{Context assumptions.}
\label{tab:context}
\centering
\small
\begin{tabularx}{\columnwidth}{@{} c >{\hsize=0.80\hsize\raggedright\arraybackslash}X >{\hsize=1.20\hsize\raggedright\arraybackslash}X @{}}
\toprule
\textbf{ID} & \textbf{Assumption} & \textbf{Value} \\
\midrule
A1 & Reference vehicle     & Mid-segment ICE passenger car, German market \\
A2 & Annual mileage        & 13{,}140 km/yr~\cite{dat2026report} \\
A3 & First-owner horizon   & 8 years~\cite{mckinsey2025dealer} \\
A4 & Workshop visits       & 1.6/yr~\cite{dat2026report} \\
A5 & Considered wear parts & Brake pads/discs, tyres, wipers, 12\,V battery, spark plugs \\
A6 & Service gross margin  & 45\,\% (published range 40--55\,\%; lower bound used as conservative estimate for mid-segment, non-premium dealers~\cite{mckinsey2025dealer,zdkzahlen2024}) \\
A7 & Franchised share      & 47\,\%~\cite{dat2026report} \\
\bottomrule
\end{tabularx}
\end{table}
 
The wear parts in A5 are deliberately limited to components whose wear is
gradual, usage-driven and not readily perceived by the driver --- the class for
which a forecast adds the most decision value and which dominates both
between-interval failures and the most common defect categories in the mandatory
German vehicle safety inspection (Hauptuntersuchung, HU: brakes, lighting,
tyres)~\cite{tuev2026report}. Prescriptive maintenance and powertrain overhauls
are out of scope. Typical replacement intervals and representative German market
prices (parts plus labour, incl.\ VAT) are given in Table~\ref{tab:wear} for the
reference vehicle (A1). All values are midpoints of published German market
ranges~\cite{adac2024repair}; actual intervals and prices vary substantially by
vehicle class---compact and entry-level vehicles tend toward the lower end of each
range, while premium and SUV segments tend toward the upper end. Where a published
range exists, a representative midpoint value is used; the sensitivity analysis
brackets this residual uncertainty.
 
\begin{table}[ht!]
\caption{Wear-part assumptions for reference vehicle A1 (mid-segment ICE). All values are midpoints of published German market ranges~\cite{adac2024repair}; actual intervals and prices vary by vehicle class}
\label{tab:wear}
\centering
\begin{tabularx}{\columnwidth}{@{} >{\raggedright\arraybackslash}X c c r @{}}
\toprule
\textbf{Part (A5)} & \textbf{Interval} & \textbf{Events/yr} (at A2) & \textbf{Price} \\
\midrule
Brake pads (per axle)      & 40{,}000 km  & 0.33 & \eur\,250 \\
Brake discs (+pads, front) & 120{,}000 km & 0.11 & \eur\,450 \\
Tyres (set of 4)           & 45{,}000 km  & 0.29 & \eur\,600 \\
Wiper blades (pair)        & 1 year       & 1.00 & \eur\,25  \\
12\,V battery (fitted)     & 6 years      & 0.17 & \eur\,150 \\
Spark plugs (set)          & 60{,}000 km  & 0.22 & \eur\,150 \\
\bottomrule
\end{tabularx}
\end{table}
 
For the mechanism, we model a single lever: the \emph{repair capture rate}~$c$,
the share of addressable wear-part work performed at the franchised dealership
rather than lost to an independent garage. Its complement $(1-c)$ is the
leakage. The current capture rate, the forecast accuracy, and the proactive-offer
conversion are stated in Table~\ref{tab:mech}; the target capture rate is not
assumed independently but derived from them.

The target capture rate is derived, not assumed: of the leaked share
$(1-\cIst)$, the fraction that PdM forecasts ($a$) and that converts on a
proactive offer ($k$) is recaptured. We set $a{=}1$ because customers cannot
distinguish a true-positive forecast from a false-positive recommendation---every
offer is evaluated and accepted or declined at rate $a \cdot k$ regardless of
whether the underlying wear has fully materialised:
\begin{align}
\cSoll &= \cIst + (1-\cIst)\cdot a \cdot k \notag \\
       &= 0.40 + 0.60 \cdot 1.0 \cdot 0.50 \notag \\
       &= 0.40 + 0.30 = 0.70 .
\end{align}
\vspace{-\intextsep}
\begin{table}[h!]
\caption{Mechanism assumptions (Current State $\rightarrow$ Target State).}
\label{tab:mech}
\centering
\small
\begin{tabularx}{\columnwidth}{@{} c >{\raggedright\arraybackslash}X c >{\raggedright\arraybackslash}p{0.27\columnwidth} @{}}
\toprule
\textbf{ID} & \textbf{Assumption (Symbol)} & \textbf{Value} & \textbf{Basis} \\
\midrule
A8  & Current capture rate ($\cIst$)       & 0.40            & ${<}$47\,\% share~\cite{dat2026report} \\
A9  & PdM forecast accuracy ($a$)          & 1.0 (model)     & reported 89--92\,\%~\cite{abougrad2026xai}; set to 1 (see text) \\
A10 & Proactive-offer conversion ($k$)     & 0.50            & Conserv.\ est.~\cite{challagalla_proactive_2009,gonzalez_service_2015} \\
--- & Target capture rate ($\cSoll$, deriv.)& $0.70$         & $\cIst+(1{-}\cIst)\cdot a\cdot k$,\;$a{=}1$ \\
\bottomrule
\end{tabularx}
\end{table}

This construction keeps the estimate conservative and never approaches full capture.
The conversion rate ($k = 0.50$, A10) is a conservative estimate: proactive,
personalised maintenance contact from the customer's own dealership is consistent
with findings that proactive contact aids uptake and loyalty~\cite{challagalla_proactive_2009,gonzalez_service_2015},
yet price competition from independent garages and customer inertia toward switching service providers prevent full conversion, making 0.50 a deliberate floor rather than a central estimate. The value of $k$ is not empirically measured in this context; calibrating it against real customer-behavior data is identified as a direct future-work priority (Section~\ref{sec:future}).

\subsection{Potential Assessment}

The calculation proceeds in three steps: the addressable wear-part value per
vehicle, the captured value under each state, and the resulting KPIs at vehicle,
fleet and lifecycle level.
 
\paragraph*{Step 1 --- Addressable wear-part value}
For each part $i$, the expected annual revenue is its replacement rate
$\lambda_i\;[\mathrm{yr}^{-1}]$ multiplied by its unit price $p_i$; both are
from Table~\ref{tab:wear}, with $\lambda_i$ computed on unrounded mileage~(A2).
Summing across~A5 yields the annual addressable wear-part revenue per vehicle:
\begin{gather*}
V_{\mathrm{wear}} = \sum_{i}\,\lambda_i \cdot p_i \\
= 82.1 + 49.3 + 175.2 + 25.0 + 25.0 + 32.9 \\
\approx \eur\,389 \quad \text{per vehicle per year.}
\end{gather*}
 
\paragraph*{Step 2 --- Captured value}
Applying the capture rates from Table~\ref{tab:mech}:
\begin{align}
\text{Captured}_{\mathrm{cur}} &= \cIst  \cdot V_{\mathrm{wear}} = 0.40 \cdot 389 \approx \eur\,156~\text{/veh/yr,} \\
\text{Captured}_{\mathrm{tgt}} &= \cSoll \cdot V_{\mathrm{wear}} = 0.70 \cdot 389 \approx \eur\,273~\text{/veh/yr.}
\end{align}
Applying the service gross margin (A6) gives the captured gross profit.
 
\paragraph*{Step 3 --- KPIs}
Table~\ref{tab:kpi} summarises the comparison. Revenue and gross-profit figures
are per vehicle per year; the fleet column scales the gross-profit delta
linearly to a representative franchised group of 5{,}000 active customer
vehicles.
 
\begin{table}[h!]
\caption{Current-state / target-state KPI comparison. Delta figures computed directly on unrounded intermediate values; current and target cells rounded independently.}
\label{tab:kpi}
\centering
\small
\begin{tabularx}{\columnwidth}{@{} >{\raggedright\arraybackslash}X c c c @{}}
\toprule
\textbf{KPI} & \textbf{Current} & \textbf{Target} & \textbf{$\Delta$} \\
\midrule
Repair capture rate              & 40\,\%     & 70\,\%      & +30\,pp \\
Wear-part revenue (/veh/yr)      & \eur\,156  & \eur\,273   & +\eur\,117 \\
Gross profit (/veh/yr)           & \eur\,70   & \eur\,123   & +\eur\,53 \\
Fleet GP, 5{,}000 veh (/yr)      & \eur\,351k & \eur\,613k  & +\eur\,263k \\
Customer touchpoints (/yr)       & 1.6        & $\sim$2.7   & +1.1 \\
First-pass inspection rate       & baseline   & improved    & qual. \\
Lifecycle retention~\cite{mckinsey2025dealer} & $\sim$25\,\% & sustained & qual. \\
\bottomrule
\end{tabularx}
\end{table}
 
The proactive touchpoint count rises because each high-value forecast event
(the wear parts in Table~\ref{tab:wear} excluding the wiper case,
$\approx$~1.1 events/yr) triggers an outreach contact that did not exist in the
interval-based current state, increasing the meaningful contact frequency
from 1.6 to about 2.7 per year, assuming forecast contacts fall outside
scheduled service visits; overlap would reduce the net gain modestly. Each additional contact is also an advisory and
upselling opportunity; the parts targeted---brakes and tyres above all---are
among the most common defect categories in the German mandatory vehicle safety
inspection (HU)~\cite{tuev2026report}, so proactive replacement directly improves the
customer's first-pass inspection rate, a tangible and communicable benefit---currently
60\,\% of franchised-dealer customers report not being approached about any
additional service at their last visit~\cite{dat2026report}.
 
\paragraph*{Sensitivity}
The result depends most on the target capture rate. Varying $\cSoll$ across a
plausible band while holding $V_{\mathrm{wear}}$ and the 45\,\% margin (A6) fixed
gives a per-vehicle gross-profit delta of:
\begin{align*}
\cSoll = 0.50 &\;\rightarrow\; \Delta \text{ gross profit} \approx \eur\,18~\text{/veh/yr,} \\
\cSoll = 0.60 &\;\rightarrow\; \Delta \text{ gross profit} \approx \eur\,35~\text{/veh/yr,} \\
\cSoll = 0.70 &\;\rightarrow\; \Delta \text{ gross profit} \approx \eur\,53~\text{/veh/yr,} \\
\cSoll = 0.80 &\;\rightarrow\; \Delta \text{ gross profit} \approx \eur\,70~\text{/veh/yr.}
\end{align*}
Even at the conservative end the dealership captures additional gross profit on
the order of tens of euros per vehicle per year, scaling to a six-figure annual
sum across a mid-size fleet---compounding effects are likely given that 90\,\% of
German vehicle owners consistently return to the same workshop~\cite{dat2026report}. The
economic case is therefore robust to the assumptions rather than an artefact of
any single optimistic value.

\subsection{Limitations}\label{sec:lim}

The analysis is intentionally assumption-driven and carries corresponding
limitations. First, the figures are illustrative: $V_{\mathrm{wear}}$, the capture
rates and the conversion factor are anchored to published market and accuracy
data but are not fitted to an operational dataset, and no prototype was built;
the contribution is a transparent model of the potential, not a measured
outcome. The first-owner horizon (A3: 8~years) is an upper-bound lifecycle horizon; the
DAT Report~2026 records a median holding period of 6.2~years~\cite{dat2026report},
so fleet-lifetime totals are correspondingly optimistic. Second, the PdM forecast accuracy (A9, 90\,\%) is reported for transparency
but does not enter the capture-rate formula. Because a proactive service
recommendation looks identical to the customer whether the underlying forecast
is correct or slightly premature, both true-positive and false-positive
predictions generate the same type of offer, accepted or declined at rate $k$.
Accuracy will improve continuously as in-service data accumulates~\cite{abougrad2026xai},
but the model outcome does not depend on reaching any particular accuracy threshold. A related long-run risk is trust erosion: if the system persistently recommends replacements the customer later judges premature, the effective conversion rate $k$ will fall below the modelled floor over repeated interactions; the sensitivity analysis brackets this effect indirectly through the lower-$\cSoll$ scenarios. Third, and most
consequential for practice, the target state presumes the dealership can actually
obtain the vehicle's telemetry. Under the EU Data Act (applicable from
12~September~2025) the user is the vehicle owner or driver, who may authorise a
third party such as the dealership to receive in-vehicle data on fair,
reasonable and non-discriminatory terms, while the OEM remains the data holder;
where the data is personal, the GDPR additionally requires a lawful basis~\cite{eu2023dataact}.
Access is thus contingent on customer consent, OEM provisioning in a usable
format, and business-to-business (B2B) compensation --- a non-technical adoption barrier that the model
does not price in. Finally, two scope caveats apply: tyres, though included as a
wear part, are often served by specialist channels and may leak independently of
the dealership's PdM capability; the shift to BEVs will erode ICE-specific wear-part volume and shift the addressable basket toward brakes, tyres, and the high-voltage battery.
\FloatBarrier
% ============================================================
\section{Discussion and Recommendations}
% ============================================================

\subsection{Interpretation of Results}

The per-vehicle gross-profit uplift of \euro\,53 per year scales to \euro\,263k
annually across a 5{,}000-vehicle fleet and compounds to approximately \euro\,2.1
million over the eight-year first-owner horizon (A3). The parts margin is,
however, only the visible component of the value: each of the additional 1.1
contact events per vehicle per year is simultaneously an advisory opportunity and
a loyalty-reinforcing interaction, consistently linked to increased repurchase
intent at the originating dealership~\cite{challagalla_proactive_2009,gonzalez_service_2015}.
The sensitivity analysis confirms robustness: even at $\cSoll = 0.50$, the
fleet-level gain remains approximately \euro\,90k per year---a positive business
case under any plausible conversion scenario.

\subsection{The Non-Technical Barrier}

The DT/ML/OTA technology stack is commercially mature and not the binding
constraint~\cite{lopes2025vehicledt,abougrad2026xai}. The real bottleneck is
telemetry access: under the EU Data Act (applicable from 12~September 2025),
in-vehicle data reaches the dealer only with the owner's explicit consent and on
OEM-provisioned terms~\cite{eu2023dataact}, making the dealer dependent on two
actors it does not control. The practical recommendation is therefore to build
consent infrastructure---opt-in flows integrated into the CRM, clear customer
value propositions---as a prerequisite rather than an afterthought, independently
of any specific DT deployment decision.

\subsection{Implications for Practice and Research}

Practitioners should position DT-based PdM as a CRM extension rather than an IT
project. The measurable value driver is the quality of proactive outreach, not
forecast accuracy~\cite{challagalla_proactive_2009,gonzalez-pridaDigitalTransformationAftersales2025};
appointment-scheduling integration and communication design should therefore
precede algorithm tuning. A phased entry via OEM brands that already expose
connected-car telemetry APIs allows conversion rates to be measured before a
broader infrastructure commitment is made.

For researchers, the priority is empirical validation of the conversion
rate~$k$, which is modelled here but not fitted to operational data. A field
study measuring opt-in and outreach acceptance rates would directly calibrate
this parameter. A second priority is extending the framework to battery electric
vehicles, where the high-voltage battery's state of health replaces spark plugs
as the dominant prognostic target.


% ============================================================
\section{Conclusion \& Future Work}
% ============================================================

\subsection{Conclusion}

This paper examined the effect of digital-twin-based predictive maintenance on
automotive aftersales, focusing on the German passenger-car market. A lifecycle
comparison showed that the transition from interval-based servicing to a
DT/PdM-enabled target state redirects wear-part work from the independent
aftermarket back to the franchised dealership by replacing reactive,
fault-driven appointments with proactive, evidence-based service recommendations.
A quantitative potential assessment, anchored to published German market data,
translated this shift into concrete KPIs: under a conservative conversion
assumption ($k = 0.50$), the repair capture rate rises from 40\,\% to 70\,\%,
increasing gross profit by approximately \euro\,53 per vehicle per year and
\euro\,263k annually across a 5\,000-vehicle fleet. Customer touchpoints grow
from 1.6 to roughly 2.7 per year, each an additional advisory and upselling
opportunity. The sensitivity analysis confirmed the economic case is robust
across the full plausible parameter range. The dominant adoption barrier is
non-technical: access to in-vehicle telemetry requires customer consent and OEM
provisioning under the EU Data Act and GDPR, a hurdle any real-world deployment
must resolve. In answer to the research question: the effect of DT-based PdM on
automotive aftersales is a measurable shift of wear-part revenue from the
independent aftermarket to the franchised dealership, accompanied by a structural
increase in customer contact frequency that amplifies retention and upselling value
beyond the parts margin alone.

\subsection{Future Work}\label{sec:future}

Three directions stand out for future research. First, the assumption framework
should be validated against an operational DT deployment, replacing modelled
KPIs with measured outcomes and calibrating the conversion rate $k$ under real
customer-behavior conditions. Second, the data-access challenge warrants
dedicated study: viable consent architectures and OEM-dealer data-sharing
agreements under the EU Data Act need to be developed and their effect on
customer opt-in quantified. Third, the framework should be extended to battery
electric vehicles, where the high-voltage battery's state of health becomes the
dominant prognostic target and the addressable wear-part basket shifts away from
ICE-specific components such as spark plugs.


% ============================================================
%\section*{Acknowledgment}
% ============================================================
%The authors would like to thank Eduard Dobenko (Mercedes-Benz) for supervising this work.

% Quellen in references.bib eintragen, hier nichts aendern
\bibliographystyle{IEEEtran}
\bibliography{references}
 
\onecolumn
\newpage
\section*{Verwendung Generativer KI-Systeme}
 
Bei der Erstellung der eingereichten Arbeit haben wir die nachfolgend aufgeführten
auf künstlicher Intelligenz (KI) basierten Systeme benutzt:
 
\begin{enumerate}
  \item Claude (Anthropic)
\end{enumerate}
 
Wir erklären, dass wir
\begin{itemize}
  \item uns aktiv über die Leistungsfähigkeit und Beschränkungen der oben genannten
        KI-Systeme informiert haben,
  \item die aus den oben angegebenen KI-Systemen direkt oder sinngemäß übernommenen
        Passagen gekennzeichnet haben,
  \item überprüft haben, dass die mithilfe der oben genannten KI-Systeme generierten
        und von uns übernommenen Inhalte faktisch richtig sind,
  \item uns bewusst sind, dass wir als Autorinnen bzw. Autoren dieser Arbeit die
        Verantwortung für die in ihr gemachten Angaben und Aussagen tragen.
\end{itemize}
 
Die oben genannten KI-Systeme haben wir wie im Folgenden dargestellt eingesetzt:
 
\vspace{0.5em}
\noindent\textbf{Tabelle~I. Einsatz von KI-Systemen}
 
\vspace{0.3em}
\noindent
\begin{tabular}{|p{0.28\textwidth}|p{0.28\textwidth}|p{0.35\textwidth}|}
\hline
\textbf{Arbeitsschritt} & \textbf{Eingesetzte(s) KI-System(e)} & \textbf{Beschreibung der Verwendungsweise} \\
\hline
Formatierung und Satz & Claude (Anthropic) & Umwandlung von Fließtext in \LaTeX{} (Tabellen, Formeln, Literaturverzeichnis) \\[4ex]
\hline
Verifikation der Berechnungen & Claude (Anthropic) & Prüfung der Zwischenrechnungen ($V_{\mathrm{wear}}$, Capture-Raten) auf Konsistenz \\[4ex]
\hline
Sprachliche Überarbeitung & Claude (Anthropic) & Korrektur von Grammatik und Stil einzelner Abschnitte \\[4ex]
\hline
\end{tabular}
 
\end{document}